\section{Related Work}
\label{sec:related}

\paragraph{Realistic input values.} Closest to us is work that makes REST
inputs look like real data. ARTE~\cite{arte} finds realistic values for
parameters by querying a public knowledge base with the parameter's name and
description. NLP2REST~\cite{nlp2rest} reads the natural-language descriptions
in a specification and turns them into extra rules, such as value ranges and
formats, that tools can use. We share their goal of getting past input
validation, but \tool uses only what the schema states (formats, patterns,
limits, examples) and the parameter names, with no network access and no
language model. This fits the REST League sandbox, which allows no remote
services, and it costs microseconds per request. The price is that \tool
cannot invent a meaningful free-text value, such as a real city for a
parameter called \texttt{location}.

\paragraph{Operation dependencies.} RESTler~\cite{restler} infers
producer--consumer relations from the specification and uses the ids in
responses to build sequences of requests. RestTestGen~\cite{resttestgen}
builds an operation dependency graph from matching field names and tests
operations in its order. Morest~\cite{morest} keeps a property graph that it
updates with what it learns during the run. ARAT-RL~\cite{arat-rl} uses
reinforcement learning to decide which operations and parameters to try and
reuses values seen in responses. Martin-Lopez et al.\ go one level deeper and
model the constraints between parameters of the same
operation~\cite{idl}. We use a deliberately small version of
producer--consumer order. What we add is a measurement of how much order
contributes once ids are already reused: on the APIs we tested, less than one
operation (Observation~\ref{obs:order}).

\paragraph{Complete testing tools.} EvoMaster~\cite{evomaster} generates
tests with an evolutionary search guided by code coverage, which needs access
to the code.
Schemathesis~\cite{schemathesis} derives property-based tests from the
schema. RESTest~\cite{restest} generates tests from a specification extended
with inter-parameter constraints. DeepREST~\cite{deeprest} learns a testing
policy with deep reinforcement learning. RestTestGen separates nominal and
error testing~\cite{resttestgen-stvr}, which our \mode{VALID} and
\mode{INVALID} modes follow. All of these are complete tools with their own
input generators. \tool is one component of a tool, and we evaluate it on
its own, with the other parts held fixed, so that its effect is not mixed
with that of a search strategy.

\paragraph{Empirical comparisons.} Corradini et al.\ compared black-box
REST testing tools on the same APIs~\cite{corradini-comparison}. Kim et
al.\ compared tools on real-world services and listed open problems for the
field, among them producing valid input values and handling dependencies
between operations~\cite{no-time-to-rest}. RESTgym~\cite{restgym} and REST
League~\cite{restleague} provide the shared infrastructure and scoring that
we follow. Their studies compare complete tools; we compare configurations of
one component on the same runner, which shows which part of the input side
moves the score.
